\section{Diagnosis: Where the Time Actually Goes}\label{sec:diagnosis}

\subsection{Methodology}

We instrumented \cref{alg:par} to attribute wall-clock time to each phase
(expand, frontier gather, node compaction, zone reclamation) and, inside the
parallel region, to each thread.  Per-thread \emph{busy} time is accumulated per
loop iteration rather than after the OpenMP barrier, so it measures work done
rather than time until the barrier; the gap between $\max_t \mathrm{busy}_t$ and
the phase wall-clock is idle time attributable to imbalance.  Inside
$\mathrm{add}$ we separately time unpacking, the three folds, and interning, and
count calls, early rejections, and antichain lengths.  All measurements are
median-of-three on an Apple M4~Pro (8 performance + 4 efficiency cores, 12
threads), \texttt{-O3}, LLVM~OpenMP.

We report arbiter15 and arbiter17 (\num{278528} and \num{1245184} symbolic
states).  Both are diagonal-free, so the $\preceq_{aLU}$ path is active.  As a
correctness cross-check throughout, UPPAAL~5.0 stores \emph{exactly} the same
number of symbolic states on these models, so the two tools are doing the same
amount of exploration and per-state cost is directly comparable.

\subsection{Finding 0: scheduling is not the bottleneck}

\begin{table}[t]
\centering
\caption{Expansion-phase parallelism, before any optimisation.  ``Imbalance
loss'' is $\max_t \mathrm{busy}_t - \tfrac{1}{T}\sum_t \mathrm{busy}_t$.}
\label{tab:balance}
\begin{tabular}{lrrrr}
\toprule
model & expand wall & $\sum_t \mathrm{busy}_t$ & parallel eff. & imbalance loss \\
\midrule
arbiter15 & \SI{1.14}{\second} & \SI{13.65}{\second} & $99.5\%$ & \SI{2}{\milli\second} \\
arbiter17 & \SI{6.95}{\second} & \SI{83.24}{\second} & $99.8\%$ & \SI{4}{\milli\second} \\
\bottomrule
\end{tabular}
\end{table}

\Cref{tab:balance} settles the question the introduction posed.  Threads are
busy $99.5$--$99.8\%$ of the expansion phase, and cumulative idle time from
uneven chunk assignment is single-digit milliseconds in multi-second runs.
Varying the \texttt{schedule(dynamic,$c$)} chunk size over
$c \in \{1,2,4,8,16,32,64\}$ changes end-to-end runtime by less than $3\%$
(\SIrange{1.29}{1.33}{\second} on arbiter15), which is the same conclusion
reached independently: the work-distribution mechanism is not what limits this
engine.

We stress this because it is the hypothesis a practitioner reaches for first,
and because it is cheap to ``fix'' something that is not broken.  Dynamic
scheduling is already doing its job.

\subsection{Finding 1: a latched reclamation trigger}

Phase attribution then shows an unmistakable serial term:

\begin{center}
\begin{tabular}{lrrr}
\toprule
model & expand & zone reclamation & serial share \\
\midrule
arbiter15 & \SI{1.14}{\second} & \SI{0.26}{\second} & $19.5\%$ \\
arbiter17 & \SI{6.95}{\second} & \SI{2.82}{\second} & $29.8\%$ \\
\bottomrule
\end{tabular}
\end{center}

The mark-and-sweep over the zone-intern table was guarded by
\texttt{interned\_count > 50000}, where \texttt{interned\_count} is a
\emph{high-water} index counter that never decreases (reclamation recycles
indices into a free list rather than renumbering).  The guard is therefore
\emph{latched}: it becomes true once, on an early level, and remains true for
the rest of the run.  A fully serial pass over all antichain nodes plus a
rebuild of the intern hash table then executes on every level.

Its cost is $O(\text{arena} + \text{table})$, which grows with the state space,
while the expansion phase it competes with is parallel.  That is why the
symptom worsens with model size---aggregate CPU utilisation falls from $835\%$
at arbiter15 to $741\%$ at arbiter17---and why it is invisible on the small
models one uses while developing.

\subsection{Finding 2: the compression tax}\label{sec:tax}

With reclamation fixed (\cref{sec:gc}), $\mathrm{add}$ accounts for $70$--$72\%$
of all CPU time in the expansion phase, against $18\%$ for successor generation
itself---that is, the passed-list bookkeeping costs four times the actual model
checking.  Instrumenting inside $\mathrm{add}$ explains why.

\begin{table}[t]
\centering
\caption{Inside $\mathrm{add}$ on arbiter17, before the fast path.  Times are
summed over all 12 threads.}
\label{tab:addprof}
\begin{tabular}{lr}
\toprule
$\mathrm{add}$ calls & \num{12255232} \\
\quad new discrete state & \num{1245183} ($10.2\%$) \\
\quad rejected by fold~(a) & \num{11010049} ($89.8\%$) \\
mean antichain length & $1.00$ \\
\midrule
time unpacking stored zones & \SI{36.33}{\second} \\
time in folds (a)--(c) & \SI{1.64}{\second} \\
time interning survivors & \SI{11.30}{\second} \\
\bottomrule
\end{tabular}
\end{table}

\Cref{tab:addprof} contains the paper's central measurement.  Three facts
combine badly:

\begin{enumerate}
  \item \textbf{Antichains are trivial.}  Mean length $1.00$: convex-union
        merging collapses each discrete state's zone set to a single zone.  The
        folds have nothing to iterate.
  \item \textbf{The common case is rejection.}  $89.8\%$ of insertions are
        duplicates killed by the first test in fold~(a).
  \item \textbf{Reaching that test requires a closure.}  Each probe unpacks the
        stored zone: scatter $|G|$ constraints, then Floyd--Warshall at
        $O(n^3)$, $n = 17$.
\end{enumerate}

The result is \SI{36.33}{\second} of decompression feeding \SI{1.64}{\second} of
comparison---a factor of $22$.  Compression bought a $12\times$ memory
reduction and charged an $O(n^3)$ toll on every one of \num{12.3}{M} queries.
We call this the \emph{compression tax}: the representation is evaluated on the
axis it optimises (bytes resident) and silently taxed on an axis nobody
measures (cost per query).

Parallelism makes the tax harder to notice, not easier.  The engine still
scales at $99.8\%$ efficiency; it simply scales a constant factor that should
not be there.  Comparing against UPPAAL makes it concrete: on arbiter15 UPPAAL
uses \SI{2.06}{\second} of CPU to store the same \num{278528} states, where our
expansion phase alone burned \SI{13.65}{\second} across twelve threads.  We were
spending roughly $6.6\times$ more CPU per state and recovering it with cores.

\subsection{A hypothesis that failed}

For the record, our first response to \cref{tab:addprof} was to eliminate
\emph{repeated} unpacking: fold~(a), (b) and (c) each walked the antichain, and
(c) restarted after every absorption, so we restructured $\mathrm{add}$ to
unpack each node once into a reusable array and fold over that.  It changed
nothing measurable (\SIrange{53.7}{56.5}{\second} in $\mathrm{add}$, within
noise).  The reason is item~1 above: with mean antichain length $1.00$ there is
no repetition to remove.  We keep the restructuring---it is strictly better when
antichains are long, as on merge-heavy models---but it is not the fix.

We also simulated a per-thread direct-mapped cache from zone index to unpacked
matrix.  Hit rates were $2.1\%$, $4.6\%$, $11.0\%$ and $21.6\%$ at $1$K, $4$K,
$16$K and $64$K entries per thread; even the last, which would cost about a
gigabyte of resident cache across twelve threads, leaves four fifths of the
probes paying full price.  Locality across the frontier is simply too poor.  The
tax has to be removed, not amortised.
